% Part: first-order-logic
% Chapter: syntax-and-semantics
% Section: subformulas

\documentclass[../../../include/open-logic-section]{subfiles}

\begin{document}

\olfileid[hi]{fol}{syn}{sbf}

\olsection{\printtoken{P}{subformula} की धारणा}

\begin{explain}
किसी दिए हुए !!{formula} को ``बनाने वाले'' !!{formula}s की चर्चा करना
प्रायः उपयोगी होता है। इन्हें हम उसके \emph{!!{subformula}} कहते हैं। हर
!!{formula} स्वयं अपना !!a{subformula} माना जाता है; $!A$ का ऐसा उपसूत्र जो
स्वयं $!A$ न हो, \emph{उचित !!{subformula}} कहलाता है।
\end{explain}

\begin{defn}[तात्कालिक !!^{subformula}]
यदि $!A$ !!a{formula} है, तो $!A$ के \emph{तात्कालिक !!{subformula}}
निम्न प्रकार आगमनात्मक रूप से परिभाषित हैं:
\begin{enumerate}
\item परमाण्विक !!{formula}s के कोई तात्कालिक !!{subformula} नहीं होते।

\tagitem{prvNot}{\indcase{!A}{\lnot !B}{$\indfrm$ का एकमात्र तात्कालिक
    !!{subformula}~$!B$ है।}}{}

\item \indcase{!A}{(!B \ast !C)}{$\indfrm$ के तात्कालिक
  !!{subformula} $!B$ और $!C$ हैं ($\ast$ दो-स्थानी संयोजकों में से कोई
  एक है)।}

\tagitem{prvAll}{\indcase{!A}{\lforall[x][!B]}{$\indfrm$ का एकमात्र
    तात्कालिक !!{subformula}~$!B$ है।}}{}

\tagitem{prvEx}{\indcase{!A}{\lexists[x][!B]}{$\indfrm$ का एकमात्र
    तात्कालिक !!{subformula}~$!B$ है।}}{}
\end{enumerate}
\end{defn}

\begin{defn}[उचित !!^{subformula}]
यदि $!A$ !!a{formula} है, तो $!A$ के \emph{उचित !!{subformula}}
निम्न प्रकार पुनरावर्ती रूप से परिभाषित हैं:
\begin{enumerate}
\item परमाण्विक !!{formula}s के कोई उचित !!{subformula} नहीं होते।

\tagitem{prvNot}{\indcase{!A}{\lnot !B}{$\indfrm$ के उचित
    !!{subformula}s में~$!B$ और~$!B$ के सभी उचित !!{subformula} आते हैं।}}{}

\item \indcase{!A}{(!B \ast !C)}{$\indfrm$ के उचित !!{subformula}s में
  $!B$, $!C$, तथा $!B$ और~$!C$ के सभी उचित !!{subformula} आते हैं।}

\tagitem{prvAll}{\indcase{!A}{\lforall[x][!B]}{$\indfrm$ के उचित
    !!{subformula}s में~$!B$ और~$!B$ के सभी उचित !!{subformula} आते हैं।}}{}

\tagitem{prvEx}{\indcase{!A}{\lexists[x][!B]}{$\indfrm$ के उचित
    !!{subformula}s में~$!B$ और~$!B$ के सभी उचित !!{subformula} आते हैं।}}{}
\end{enumerate}
\end{defn}

\begin{defn}[!!^{subformula}]
$!A$ के !!{subformula}s में स्वयं $!A$ और उसके सभी उचित !!{subformula} आते
हैं।
\end{defn}

\begin{explain}
तात्कालिक !!{subformula}s और उचित !!{subformula}s की परिभाषाओं में सूक्ष्म
अंतर पर ध्यान दें। पहली स्थिति में हमने~$!A$ के प्रत्येक संभव रूप के लिए
सूत्र~$!A$ के तात्कालिक !!{subformula} सीधे परिभाषित किए हैं। यह मामलों
द्वारा दी गई स्पष्ट परिभाषा है, और ये मामले !!{formula}s के समुच्चय की
आगमनात्मक परिभाषा का अनुसरण करते हैं। दूसरी स्थिति में भी हमने सभी
!!{formula}s के समुच्चय की परिभाषा का अनुसरण किया है, किंतु हर मामले में
छोटे !!{formula}s $!B$, $!C$ के साथ-साथ इन !!{formula}s के उचित
!!{subformula} भी सम्मिलित किए हैं। इससे परिभाषा \emph{पुनरावर्ती} बनती है। सामान्यतः किसी
आगमनात्मक रूप से परिभाषित समुच्चय (हमारे मामले में !!{formula}s) पर फलन की
परिभाषा तब पुनरावर्ती होती है जब फलन की परिभाषा के मामले स्वयं उसी फलन का
उपयोग करें। किंतु परिभाषा सुपरिभाषित होने के लिए यह सुनिश्चित करना आवश्यक है
कि हम फलन के मान केवल उन तर्कों के लिए उपयोग करें जो परिभाषित किए जा रहे तर्क
से ``पहले'' आते हैं---हमारे मामले में $(!B \ast !C)$ के लिए ``उचित
!!{subformula}'' परिभाषित करते समय हम केवल ``पहले के'' !!{formula}s $!B$ और
$!C$ के उचित !!{subformula}s का उपयोग करते हैं।
\end{explain}

\begin{prop}
\ollabel{prop:subfrm-trans}
मान लें कि $!B$, $!A$ का उपसूत्र है और $!C$, $!B$ का उपसूत्र है। तब $!C$,
$!A$ का उपसूत्र है। दूसरे शब्दों में, उपसूत्र-संबंध संक्रामक है।
\end{prop}

\begin{prob}
\olref[fol][syn][sbf]{prop:subfrm-trans} सिद्ध कीजिए।
\end{prob}

\begin{prop}
\ollabel{prop:count-subfrms}
मान लें कि $!A$ ऐसा सूत्र है जिसमें $n$ संयोजक और परिमाणक हैं। तब $!A$ के
अधिक-से-अधिक $2n+1$ उपसूत्र हैं।
\end{prop}

\begin{prob}
\olref[fol][syn][sbf]{prop:count-subfrms} सिद्ध कीजिए।
\end{prob}

\end{document}
